
est
commutatif. Soit \(\chi\in\Ghat(A)\)~; alors~: \(h\circ f_\chi=f_{\Ghat(h)(\chi)}\)~; d'où~:
\[
  H(h)(\Phi_A(\chi))=H(h)\circ H(f_\chi(\psi))
  =H(f_{\Ghat(h)(\chi)})(\psi)=\Phi_{A'}(\Ghat(h)(\chi).
\]

Il reste donc à montrer que le morphisme de foncteurs ainsi construit est compatible aux structures de groupes filtrés.

Supposons que \(\chi\in F^n\Ghat(A)\). Dire que \(\chi\) est de filtration \(\geq n\) équivaut à dire que \(f_\chi:K[X_i]_{i\geq 1}\to A\), se factorise par l'endomorphisme \(g_n\) de \(K[X_i]_{i\geq 1}\) défini par~: \(g_n(X_i)=0\) si \(1\leq i\leq n-1\), et \(g_n(X_i)=X_i\) si \(i\geq n\).
Pour montrer que \(\Phi\) est compatible aux filtrations, il suffit donc de montrer que \(H(g_n)(\Phi(\varphi))\in F^nH\). Or par construction~:
\[
  \Phi(\varphi)\equiv \Phi(1+\ldots+X_{n-1}t^{n-1})\ \mathrm{mod.}\ F^nH,
\]
d'où~:
\[
  H(g_n)(\Phi(\varphi))
  \equiv H(g_n)\bigl(\Phi(1+\ldots+X_{n-1}t^{n-1})\bigr)
  =0\ \mathrm{mod.}\ F^nH.
\]

Pour vérifier l'additivité de \(\Phi\), il suffit de le faire pour les éléments de degré fini de \(\Ghat(A)\), et de passer ensuite à la limite. Soient donc \(\chi_1,\chi_2\) deux éléments de degré \(n\) de \(\Ghat(A)\), et \(\chi\) leur produit. On considère le diagramme commutatif~:
\[
\begin{tikzcd}[column sep=tiny]
K[X_1,\ldots,X_{2n}] \arrow[rr,"f_\chi"] \arrow[d,hook] && A\\
\bigl(K[T_1,\ldots,T_{2n}]\bigr)^{\mathfrak S_n\times\mathfrak S_n}
\arrow[urr,"f'_\chi"] && ,
\end{tikzcd}
\]
où les flèches sont définies comme suit~: \(\mathfrak S_n\times\mathfrak S_n\) opère sur \(K[T_1,\ldots,T_{2n}]\) 
